% GENERATED FILE. Edit canonical lessons, then run npm run generate:books.
% canonical-source-hash: fnv1a64:04b808aa0d6eaa35
% canonical-lessons: BN-C435-ghas, BN-C435-lota, BN-C435-shakha, BN-C435-kath, BN-C435-chhata

\chapter{Grass, Creepers and Wood}
\label{ch:bn-grass-creepers-and-wood}
\hlchaptermodality{435}

\begin{chapteropening}
\textbf{By the end of this chapter:} \emph{I can talk about grass, creepers, branches and wood, and trimming a garden.}

Complete the last lesson of chapter 435: I can talk about grass, creepers, branches and wood, and trimming a garden.
\end{chapteropening}

\section[\texorpdfstring{\bn{ঘাস} (ghās)}{ঘাস (ghās)}]{\bn{ঘাস} (ghās) — grass}

\label{lesson:BN-C435-ghas}



\begin{quote}
Before the new one: say the Bengali for to spoil, then the Bengali for to conserve.
\end{quote}

\subsection*{You'll want to know: \bn{ঘাস}}
\textbf{\bn{ঘাস}} — \emph{ghās} — \textquotedblleft{}grass\textquotedblright{}.

\begin{quote}
\bn{সকালে} \bn{মাঠের} \bn{ঘাস} \bn{শিশিরে} \bn{ভেজা।}
\end{quote}

\begin{quote}
\emph{sôkāle māṭher ghās shishire bhejā.}
\end{quote}

\begin{quote}
In the morning the grass on the field is wet with dew.
\end{quote}

\begin{cousinweb}
\textbf{\bn{ঘাস}} (\emph{ghās}) is inherited: CDIAL lists Bengali \emph{ghās}, \textquotedblleft{}grass\textquotedblright{}, under Sanskrit \emph{ghāsá}, \textquotedblleft{}fodder, grass\textquotedblright{} (CDIAL 4471).
\end{cousinweb}

\begin{grammarlens}[title={What you've built}]
One more word for getting around.
\end{grammarlens}

\subsection*{Your turn}
\begin{itemize}
\raggedright
  \item \emph{Say it:} \emph{ghās}
  \item \emph{Say it:} \emph{ghās}, once more
  \item \emph{Say it:} \emph{sôkāle māṭher ghās shishire bhejā.} — In the morning the grass on the field is wet with dew.
  \item \emph{Recall:} say the Bengali for to spoil, then the Bengali for to conserve, then say \emph{ghās} again
\end{itemize}

\begin{tcolorbox}[breakable,colback=teal!4,colframe=teal!35!black,title={Before you move on}]
What does \bn{ঘাস} mean? (\textquotedblleft{}Grass\textquotedblright{}.) Say it once more.
\end{tcolorbox}
\section[\texorpdfstring{\bn{লতা} (lôtā)}{লতা (lôtā)}]{\bn{লতা} (lôtā) — a creeper; a climbing plant}

\label{lesson:BN-C435-lota}



\begin{quote}
Before the new one: say the Bengali for to conserve, then the Bengali for grass.
\end{quote}

\subsection*{You'll want to know: \bn{লতা}}
\textbf{\bn{লতা}} — \emph{lôtā} — \textquotedblleft{}a creeper; a climbing plant\textquotedblright{}.

\begin{quote}
\bn{আমাদের} \bn{বারান্দায়} \bn{একটা} \bn{ফুলের} \bn{লতা} \bn{আছে।}
\end{quote}

\begin{quote}
\emph{āmāder bārāndāy ekṭā phuler lôtā āchhe.}
\end{quote}

\begin{quote}
There is a flowering creeper on our veranda.
\end{quote}

\begin{cousinweb}
\textbf{\bn{লতা}} (\emph{lôtā}) is Sanskrit \emph{latā}, \textquotedblleft{}a creeper, any twining plant\textquotedblright{} (Monier-Williams), taken over whole.
\end{cousinweb}

\begin{grammarlens}[title={What you've built}]
One more word for getting around.
\end{grammarlens}

\subsection*{Your turn}
\begin{itemize}
\raggedright
  \item \emph{Say it:} \emph{lôtā}
  \item \emph{Say it:} \emph{lôtā}, once more
  \item \emph{Say it:} \emph{āmāder bārāndāy ekṭā phuler lôtā āchhe.} — There is a flowering creeper on our veranda.
  \item \emph{Recall:} say the Bengali for to conserve, then the Bengali for grass, then say \emph{lôtā} again
\end{itemize}

\begin{tcolorbox}[breakable,colback=teal!4,colframe=teal!35!black,title={Before you move on}]
What does \bn{লতা} mean? (\textquotedblleft{}A creeper; a climbing plant\textquotedblright{}.) Say it once more.
\end{tcolorbox}
\section[\texorpdfstring{\bn{শাখা} (shākhā)}{শাখা (shākhā)}]{\bn{শাখা} (shākhā) — a branch (of a tree, a bank or a river)}

\label{lesson:BN-C435-shakha}



\begin{quote}
Before the new one: say the Bengali for grass, then the Bengali for a creeper.
\end{quote}

\subsection*{You'll want to know: \bn{শাখা}}
\textbf{\bn{শাখা}} — \emph{shākhā} — \textquotedblleft{}a branch (of a tree, a bank or a river)\textquotedblright{}.

\begin{quote}
\bn{আমাদের} \bn{ব্যাংকের} \bn{একটা} \bn{শাখা} \bn{স্টেশনের} \bn{কাছে।}
\end{quote}

\begin{quote}
\emph{āmāder bængker ekṭā shākhā sṭeshôner kāchhe.}
\end{quote}

\begin{quote}
Our bank has a branch near the station.
\end{quote}

On a tree people mostly say \bn{ডাল}, which sounds like \bn{ডাল}, lentils; \bn{শাখা} is the careful word, and the only one for a bank or a river.

\begin{cousinweb}
\textbf{\bn{শাখা}} (\emph{shākhā}) is Sanskrit \emph{śākhā}, \textquotedblleft{}a branch\textquotedblright{} (Monier-Williams), taken over whole; a bank's branch is the same image.
\end{cousinweb}

\begin{grammarlens}[title={What you've built}]
One more word for getting around.
\end{grammarlens}

\subsection*{Your turn}
\begin{itemize}
\raggedright
  \item \emph{Say it:} \emph{shākhā}
  \item \emph{Say it:} \emph{shākhā}, once more
  \item \emph{Say it:} \emph{āmāder bængker ekṭā shākhā sṭeshôner kāchhe.} — Our bank has a branch near the station.
  \item \emph{Recall:} say the Bengali for grass, then the Bengali for a creeper, then say \emph{shākhā} again
\end{itemize}

\begin{tcolorbox}[breakable,colback=teal!4,colframe=teal!35!black,title={Before you move on}]
What does \bn{শাখা} mean? (\textquotedblleft{}A branch (of a tree, a bank or a river)\textquotedblright{}.) Say it once more.
\end{tcolorbox}
\section[\texorpdfstring{\bn{কাঠ} (kāṭh)}{কাঠ (kāṭh)}]{\bn{কাঠ} (kāṭh) — wood; timber}

\label{lesson:BN-C435-kath}



\begin{quote}
Before the new one: say the Bengali for a creeper, then the Bengali for a branch.
\end{quote}

\subsection*{You'll want to know: \bn{কাঠ}}
\textbf{\bn{কাঠ}} — \emph{kāṭh} — \textquotedblleft{}wood; timber\textquotedblright{}.

\begin{quote}
\bn{শীতে} \bn{আমরা} \bn{কাঠ} \bn{জ্বালালাম।}
\end{quote}

\begin{quote}
\emph{shīte āmrā kāṭh jvālālām.}
\end{quote}

\begin{quote}
In winter we burned wood.
\end{quote}

\bn{ঠ}, the breathy partner of \bn{ট}, closes the word.

\begin{cousinweb}
\textbf{\bn{কাঠ}} (\emph{kāṭh}) is inherited: CDIAL lists Bengali \emph{kāṭh}, \textquotedblleft{}wood\textquotedblright{}, under Sanskrit \emph{kāṣṭhá}, \textquotedblleft{}a piece of wood\textquotedblright{} (CDIAL 3120).
\end{cousinweb}

\begin{grammarlens}[title={What you've built}]
One more word for getting around.
\end{grammarlens}

\subsection*{Your turn}
\begin{itemize}
\raggedright
  \item \emph{Say it:} \emph{kāṭh}
  \item \emph{Say it:} \emph{kāṭh}, once more
  \item \emph{Say it:} \emph{shīte āmrā kāṭh jvālālām.} — In winter we burned wood.
  \item \emph{Recall:} say the Bengali for a creeper, then the Bengali for a branch, then say \emph{kāṭh} again
\end{itemize}

\begin{tcolorbox}[breakable,colback=teal!4,colframe=teal!35!black,title={Before you move on}]
What does \bn{কাঠ} mean? (\textquotedblleft{}Wood; timber\textquotedblright{}.) Say it once more.
\end{tcolorbox}
\section[\texorpdfstring{\bn{ছাঁটা} (chhā̃ṭā)}{ছাঁটা (chhā̃ṭā)}]{\bn{ছাঁটা} (chhā̃ṭā) — to trim; to prune; to cut back}

\label{lesson:BN-C435-chhata}



\begin{quote}
Before the new one: say the Bengali for a branch, then the Bengali for wood.
\end{quote}

\subsection*{You'll want to know: \bn{ছাঁটা}}
\textbf{\bn{ছাঁটা}} — \emph{chhā̃ṭā} — \textquotedblleft{}to trim; to prune; to cut back\textquotedblright{}.

\begin{quote}
\bn{মালি} \bn{বাগানের} \bn{গাছগুলো} \bn{ছাঁটলেন।}
\end{quote}

\begin{quote}
\emph{māli bāgāner gāchhgulo chhā̃ṭlen.}
\end{quote}

\begin{quote}
The gardener trimmed the plants in the garden.
\end{quote}

Past (he/she …): \bn{তিনি} \bn{গাছ} \bn{ছাঁটলেন}, \emph{tini gāchh chhā̃ṭlen}.

Also of hair and of staff: \bn{চুল} \bn{ছাঁটা}, to trim hair.

\begin{cousinweb}
\textbf{\bn{ছাঁটা}} (\emph{chhā̃ṭā}) is inherited: CDIAL lists Bengali \emph{chā̃ṭā}, \textquotedblleft{}to lop\textquotedblright{}, under the reconstructed \emph{chāṇṭ}, \textquotedblleft{}cut\textquotedblright{} (CDIAL 5013).
\end{cousinweb}

\begin{grammarlens}[title={What you've built}]
One more word for getting around.
\end{grammarlens}

\subsection*{Your turn}
\begin{itemize}
\raggedright
  \item \emph{Say it:} \emph{chhā̃ṭā}
  \item \emph{Say it:} \emph{chhā̃ṭā}, once more
  \item \emph{Say it:} \emph{māli bāgāner gāchhgulo chhā̃ṭlen.} — The gardener trimmed the plants in the garden.
  \item \emph{Recall:} say the Bengali for a branch, then the Bengali for wood, then say \emph{chhā̃ṭā} again
\end{itemize}

\begin{tcolorbox}[breakable,colback=teal!4,colframe=teal!35!black,title={Before you move on}]
What does \bn{ছাঁটা} mean? (\textquotedblleft{}To trim; to prune; to cut back\textquotedblright{}.) Say it once more.
\end{tcolorbox}
